% titel: Blatt 9, A1: Vollständige Induktion
% skript: Kap. 1.2 (Vollständige Induktion, S. 11)

% Teil I: Werkzeuge, kein Lösungsweg.

Zeigen, dass zwei Aussagen der Form ``für alle $n\in\mathbb N$ gilt …'' wahr
sind, jeweils per vollständiger Induktion.

\begin{werkzeugkasten}[Vollständige Induktion]
Eine Aussage $A(n)$ für alle $n\in\mathbb N$ beweist man in zwei Schritten:
\begin{enumerate}
  \item \textbf{Induktionsanfang} (IA): $A(1)$ direkt prüfen.
  \item \textbf{Induktionsschritt} (IS): unter der \emph{Induktionsannahme}
        $A(n)$ (bzw.\ $A(k)$) folgt $A(n+1)$ (bzw.\ $A(k+1)$).
\end{enumerate}
Antrit die Kette $A(1)\Rightarrow A(2)\Rightarrow\cdots$, so gilt $A(n)$ für
alle $n$. Oft fängt man bei $n=0$ an, wenn die Behauptung das nahelegt.

\textit{Skript:} Kap. 1.2, S. 11.
\end{werkzeugkasten}

\begin{werkzeugkasten}[Hilfsmittel im IS]
Im Schritt rechnet man die neue Summe bzw.\ das neue Produkt auf die
Induktionsannahme zurück:
\begin{itemize}
  \item \textbf{Summe:} $\displaystyle\sum_{k=1}^{n+1}a_k = a_{n+1} +
        \sum_{k=1}^{n}a_k$ — dann die Voraussetzung einsetzen.
  \item \textbf{Produkt:} $\displaystyle\prod_{k=1}^{n+1}a_k =
        a_{n+1}\cdot\prod_{k=1}^{n}a_k$ — dann die Voraussetzung einsetzen.
  \item Der letzte Summand bzw.\ Faktor hat die Form $\frac{n+1}{2^{n+1}}$
        bzw.\ $\frac{2}{n+1}$, sodass man auf einen gemeinsamen Nenner bringen
        und zusammenfassen kann.
\end{itemize}
\textit{Skript:} Kap. 1.2, S. 11--12.
\end{werkzeugkasten}

\begin{achtung}[Darauf achten]
\begin{itemize}
  \item Im Summen-Beispiel entsteht rechts ein Ausdruck
        $2-\frac{n+2}{2^n}+\frac{n+1}{2^{n+1}}$; nur mit dem gemeinsamen
        Nenner $2^{n+1}$ lassen sich die Brüche sauber zusammenfassen.
  \item Im Produkt-Beispiel schreibt man jeden Faktor um
        (hier $1+\frac2k=\frac{k+2}{k}$): so wird das Produkt zum
        Teleskopprodukt, das Schritt für Schritt kürzt.
  \item Die Behauptung ist für $n\in\mathbb N$ formuliert; prüfen, ob der
        IA bei $n=1$ beginnt.
  \item Nicht den IS als blosse Umformung der Induktionsannahme missverstehen;
        das Ziel ist die Behauptung für $n+1$.
\end{itemize}
\end{achtung}
